% page 361 of the facsimile (image p-360 of the scan)
% block 162.6 x 237.7 mm ; left margin 39.4 mm
\pgc{17.780mm}{0.000mm}{
\l{\cel{51}353}
\l{}
\l{\sou{mandato}. I. Dokumento emisita en posto-kontoro, e po qua}
\l{\cel{3}irga posto-kontoro di la lando pagos la pekunio-quanto}
\l{\cel{3}indikita sur olu. - II.Dokumento donita ad ulu, da admi-}
\l{\cel{3}nistro-fako statala, e po qua la Trezoro statala pagos}
\l{\cel{3}la pekunio-quanto indikita sur olu. -III. (\sou{banko}). Doku-}
\l{\cel{3}mento per qua bankiero imperas a sua debanto ke ica pagez}
\l{\cel{3}a ta o ca persono la pekunio-quanto indikita sur olu. -}
\l{\cel{3}DFIRS.}
\l{}
\l{\sou{mandelo}. (\sou{bot}.)Frukto di arboro ek la familio \textquotedbl{}rozacei\textquotedbl{}.}
\l{\cel{3}L.\cel{1}\sou{amygdalus communis}. DFIRS.}
\l{}
\l{\sou{mandibulo}. (\sou{anat.}) La parto infra ek la du parti osta di}
\l{\cel{3}la boko,qui suportas la denti, che la animali vertebroza.}
\l{\cel{3}- DEFIS.}
\l{}
\l{\sou{mandolino}. (\sou{muziko)}. Muzik-instrumento ek framo, konvexa sube,}
\l{\cel{3}qua havas normale mancho kun quar kordi metala duopla,}
\l{\cel{3}dispozita ed inter-akordita quale olti di violino, e quin}
\l{\cel{3}onu vibrigas per peco de ucel-plumo o di skalio. - DEFIRS.}
\l{}
\l{\sou{mandoro}. (\sou{muziko}).Muzik-instrumento (ne plus uzata) analoga a}
\l{\cel{3}mandolino. - DEFI.}
\l{}
\l{\sou{mandragoro}. (\sou{bot.}) Planto ek la genero \textquotedbl{}solanei\textquotedbl{}, kun radiko}
\l{\cel{3}pulpoza, narkotigiva e purgiva, ed a qua onu atribuis}
\l{\cel{3}efiki magiala. - L.\cel{1}\sou{mandragora officinarum}. - EFIS.}
\l{}
\l{\sou{mandrelo}. (\sou{tekn.)}\cel{2}Peco de fero fite a qua onu forjas peci}
\l{\cel{3}quin onu volas igar kava.EFIS.}
\l{}
\l{\sou{mandrino.(teknol.)}\cel{3}Peco quan onu visas sur tornilo horizon-}
\l{\cel{3}tala e per qua la axo rotacanta rotacigas la objekto quan}
\l{\cel{3}onu volas fasonar. - FIS.}
\l{}
\l{\sou{manejo}. Loko ube onu drasas la kavali, ube onu docas e prak-}
\l{\cel{3}tikas la kavalk-arto. - DEFIRS.}
\l{}
\l{\sou{manekino}. I. (1) Figuro ligna o vaxa, qua reprezentas viro o}
\l{\cel{3}muliero, olquan uzas la piktisti, la statuifisti, por probar}
\l{\cel{3}posturi, dispozar drapiraji, dum la absenteso di la homa}
\l{\cel{3}modelo. (2) Figuro quan la taliori vestizas por uzesar kom}
\l{\cel{3}modeli pri la faco di la kostumi. - II.Figuro realeso-granda,}
\l{\cel{3}sur qua onu exercas la studenti di medicino, di kirurgio,}
\l{\cel{3}pri la bandajizo o pri la movi quin postulas ula operaci. -}
\l{\cel{3}DEFRS.}
\l{}
\l{\sou{maneto}. (\sou{teknol.}) Mikra mancho quan onu sizas por movigar meka-}
\l{\cel{3}nismo. - FI.}
\l{}
\l{mango. (\sou{bot.}) Frukto di granda arboro di India e di la sudo-}
\l{\cel{3}regioni di Amerika, e qua kelke similesas persiko. -}
\l{\cel{3}L. mango. - DEFIS.}
}
